% GENERATED FILE. Edit canonical lessons, then run npm run generate:books.
% canonical-source-hash: fnv1a64:e568890c8a496b03
% canonical-lessons: KA-C163-anumanapadu, KA-C163-sahisu, KA-C163-mosa-madu, KA-C163-hondiko, KA-C163-nedu

\chapter{To doubt, To put up with, To cheat, To get along, To plant}
\label{ch:ka-to-doubt-to-put-up-with-to-cheat-to-get-along-to-plant}
\hlchaptermodality{163}

\begin{chapteropening}
\textbf{By the end of this chapter:} \emph{I can say to doubt, to put up with, to cheat, to get along and to plant.}

Complete the last lesson of chapter 163: I can say to doubt, to put up with, to cheat, to get along and to plant.
\end{chapteropening}

\section[\texorpdfstring{\kn{ಅನುಮಾನಪಡು} (anumānapaḍu)}{ಅನುಮಾನಪಡು (anumānapaḍu)}]{\kn{ಅನುಮಾನಪಡು} (anumānapaḍu) — to doubt, to suspect}

\label{lesson:KA-C163-anumanapadu}



\begin{quote}
Before the new one: say the Kannada for to take part, then the Kannada for to transfer.
\end{quote}

\subsection*{You'll want to know: \kn{ಅನುಮಾನಪಡು}}
\textbf{\kn{ಅನುಮಾನಪಡು}} — \emph{anumānapaḍu} — \textquotedblleft{}to doubt, to suspect\textquotedblright{}.

\begin{grammarlens}[title={What you've built}]
One more word for this chapter.
\end{grammarlens}

\subsection*{Your turn}
\begin{itemize}
\raggedright
  \item \emph{Say it:} \emph{anumānapaḍu}
  \item \emph{Say it:} \emph{anumānapaḍu}, once more
  \item \emph{Say it:} say \emph{anumānapaḍu}
  \item \emph{Recall:} say the Kannada for to take part, then the Kannada for to transfer, then say \emph{anumānapaḍu} again
\end{itemize}

\begin{tcolorbox}[breakable,colback=teal!4,colframe=teal!35!black,title={Before you move on}]
What does \kn{ಅನುಮಾನಪಡು} mean? (\textquotedblleft{}To doubt, to suspect\textquotedblright{}.) Say it once more.
\end{tcolorbox}
\section[\texorpdfstring{\kn{ಸಹಿಸು} (sahisu)}{ಸಹಿಸು (sahisu)}]{\kn{ಸಹಿಸು} (sahisu) — to put up with, to tolerate}

\label{lesson:KA-C163-sahisu}



\begin{quote}
Before the new one: say the Kannada for to transfer, then the Kannada for to doubt.
\end{quote}

\subsection*{You'll want to know: \kn{ಸಹಿಸು}}
\textbf{\kn{ಸಹಿಸು}} — \emph{sahisu} — \textquotedblleft{}to put up with, to tolerate\textquotedblright{}.

\begin{grammarlens}[title={What you've built}]
One more word for this chapter.
\end{grammarlens}

\subsection*{Your turn}
\begin{itemize}
\raggedright
  \item \emph{Say it:} \emph{sahisu}
  \item \emph{Say it:} \emph{sahisu}, once more
  \item \emph{Say it:} say \emph{sahisu}
  \item \emph{Recall:} say the Kannada for to transfer, then the Kannada for to doubt, then say \emph{sahisu} again
\end{itemize}

\begin{tcolorbox}[breakable,colback=teal!4,colframe=teal!35!black,title={Before you move on}]
What does \kn{ಸಹಿಸು} mean? (\textquotedblleft{}To put up with, to tolerate\textquotedblright{}.) Say it once more.
\end{tcolorbox}
\section[\texorpdfstring{\kn{ಮೋಸ} \kn{ಮಾಡು} (mōsa māḍu)}{ಮೋಸ ಮಾಡು (mōsa māḍu)}]{\kn{ಮೋಸ} \kn{ಮಾಡು} (mōsa māḍu) — to cheat}

\label{lesson:KA-C163-mosa-madu}



\begin{quote}
Before the new one: say the Kannada for to doubt, then the Kannada for to put up with.
\end{quote}

\subsection*{You'll want to know: \kn{ಮೋಸ} \kn{ಮಾಡು}}
\textbf{\kn{ಮೋಸ} \kn{ಮಾಡು}} — \emph{mōsa māḍu} — \textquotedblleft{}to cheat\textquotedblright{}.

\begin{grammarlens}[title={What you've built}]
One more word for this chapter.
\end{grammarlens}

\subsection*{Your turn}
\begin{itemize}
\raggedright
  \item \emph{Say it:} \emph{mōsa māḍu}
  \item \emph{Say it:} \emph{mōsa māḍu}, once more
  \item \emph{Say it:} say \emph{mōsa māḍu}
  \item \emph{Recall:} say the Kannada for to doubt, then the Kannada for to put up with, then say \emph{mōsa māḍu} again
\end{itemize}

\begin{tcolorbox}[breakable,colback=teal!4,colframe=teal!35!black,title={Before you move on}]
What does \kn{ಮೋಸ} \kn{ಮಾಡು} mean? (\textquotedblleft{}To cheat\textquotedblright{}.) Say it once more.
\end{tcolorbox}
\section[\texorpdfstring{\kn{ಹೊಂದಿಕೊ} (hondiko)}{ಹೊಂದಿಕೊ (hondiko)}]{\kn{ಹೊಂದಿಕೊ} (hondiko) — to get along, to adjust}

\label{lesson:KA-C163-hondiko}



\begin{quote}
Before the new one: say the Kannada for to put up with, then the Kannada for to cheat.
\end{quote}

\subsection*{You'll want to know: \kn{ಹೊಂದಿಕೊ}}
\textbf{\kn{ಹೊಂದಿಕೊ}} — \emph{hondiko} — \textquotedblleft{}to get along, to adjust\textquotedblright{}.

\begin{grammarlens}[title={What you've built}]
One more word for this chapter.
\end{grammarlens}

\subsection*{Your turn}
\begin{itemize}
\raggedright
  \item \emph{Say it:} \emph{hondiko}
  \item \emph{Say it:} \emph{hondiko}, once more
  \item \emph{Say it:} say \emph{hondiko}
  \item \emph{Recall:} say the Kannada for to put up with, then the Kannada for to cheat, then say \emph{hondiko} again
\end{itemize}

\begin{tcolorbox}[breakable,colback=teal!4,colframe=teal!35!black,title={Before you move on}]
What does \kn{ಹೊಂದಿಕೊ} mean? (\textquotedblleft{}To get along, to adjust\textquotedblright{}.) Say it once more.
\end{tcolorbox}
\section[\texorpdfstring{\kn{ನೆಡು} (neḍu)}{ನೆಡು (neḍu)}]{\kn{ನೆಡು} (neḍu) — to plant}

\label{lesson:KA-C163-nedu}



\begin{quote}
Before the new one: say the Kannada for to cheat, then the Kannada for to get along.
\end{quote}

\subsection*{You'll want to know: \kn{ನೆಡು}}
\textbf{\kn{ನೆಡು}} — \emph{neḍu} — \textquotedblleft{}to plant\textquotedblright{}.

\begin{grammarlens}[title={What you've built}]
One more word for this chapter.
\end{grammarlens}

\subsection*{Your turn}
\begin{itemize}
\raggedright
  \item \emph{Say it:} \emph{neḍu}
  \item \emph{Say it:} \emph{neḍu}, once more
  \item \emph{Say it:} say \emph{neḍu}
  \item \emph{Recall:} say the Kannada for to cheat, then the Kannada for to get along, then say \emph{neḍu} again
\end{itemize}

\begin{tcolorbox}[breakable,colback=teal!4,colframe=teal!35!black,title={Before you move on}]
What does \kn{ನೆಡು} mean? (\textquotedblleft{}To plant\textquotedblright{}.) Say it once more.
\end{tcolorbox}
